% from ega4-10-fr.tex
\oldpage[IV-3]{99}
soit $Z\neq\varnothing$ une partie fermée \emph{dans} $Y'$, et $z\in Z$~; il
y a un voisinage ouvert $U$ de $z$ dans $Y$ tel que $U\cap Y'$ soit réunion
d'un nombre fini de parties $Y'_j$ fermées dans $U$~; si $j$ est un indice tel
que $z\in Y'_j$, $Z\cap Y'_j$ est donc fermé dans $U$. Puisque $f(X)\cap U$
est très dense dans $U$ (10.1.4), $Z\cap Y'_j\cap f(X)$ n'est pas vide
(10.1.2, \emph{a)}), ni \emph{a fortiori} $Z\cap f(X)$~; mais comme
$Z\subset Y'$, on a $Z\cap f(X)=Z\cap f'(X')$~; le critère
(10.1.2, \emph{a)}) montre donc que $f'(X')$ est très dense dans $Y'$, et
l'on conclut à l'aide de (10.2.1, \emph{b)}).

\begin{corollary}[10.2.5]
\label{IV.10.2.5-fr}
Soient $f:X\to Y$ une application continue, $(V_\alpha)$ un recouvrement
ouvert de $Y$. Si, pour tout $\alpha$, la restriction
$f^{-1}(V_\alpha)\to V_\alpha$ de $f$ est un quasi-homéomorphisme, alors $f$
est un quasi-homéomorphisme.
\end{corollary}

Cela résulte aussitôt du critère (10.2.1, \emph{b)}) et de (10.1.4).

\begin{corollary}[10.2.6]
\label{IV.10.2.6-fr}
Soient $f:X\to Y$ un quasi-homéomorphisme, $Y'$ une partie localement
quasi-constructible de $Y$, $X'=f^{-1}(Y')$. Pour que $Y'$ soit quasi-compact
(resp. noethérien, resp. rétrocompact dans $Y$), il faut et il suffit que $X'$
soit quasi-compact (resp. noethérien, resp. rétrocompact dans $X$).
\end{corollary}

Démontrons d'abord les deux premières assertions~; en vertu de (10.2.4), on
peut se borner au cas où $Y'=Y$. Dire que $X$ est quasi-compact (resp.
noethérien) signifie que pour toute famille filtrante croissante $(U_\alpha)$
dans $\mathfrak O(X)$ ayant $X$ pour plus grand élément (resp. pour toute
famille filtrante croissante $(U_\alpha)$ dans $\mathfrak O(X)$), il existe
$\gamma$ tel que $U_\alpha=U_\gamma$ pour $\alpha\geq\gamma$. Comme
$U\mapsto f^{-1}(U)$ est une bijection de $\mathfrak O(Y)$ sur
$\mathfrak O(X)$, notre assertion résulte aussitôt de la remarque précédente.

Les ouverts quasi-compacts de $X$ sont donc les ensembles de la forme
$f^{-1}(U)$ où $U$ est ouvert quasi-compact dans $Y$, en vertu de
(10.2.1, \emph{a)}) et de ce qui précède. Pour que $X'$ soit rétrocompact dans
$X$, il faut et il suffit alors que pour tout ouvert quasi-compact $U$ dans
$Y$, $f^{-1}(U)\cap X'=f^{-1}(U\cap Y')$ soit quasi-compact
($\mathbf 0_{\mathrm{III}}$, 9.1.1)~; la première partie de la démonstration
montre que cela équivaut à dire que $U\cap Y'$ est quasi-compact pour tout
ouvert quasi-compact $U$, c'est-à-dire que $Y'$ est rétrocompact dans $Y$.

\begin{proposition}[10.2.7]
\label{IV.10.2.7-fr}
Soit $f:X\to Y$ un quasi-homéomorphisme. Alors l'application
$Z\mapsto f^{-1}(Z)$ de $\mathfrak P(Y)$ dans $\mathfrak P(X)$ définit par
restriction les bijections suivantes (cf. (10.1.1) pour les notations)~:
\[
\begin{gathered}
\mathfrak O(Y)\xrightarrow{\sim}\mathfrak O(X)\\
\mathfrak F(Y)\xrightarrow{\sim}\mathfrak F(X)\\
\mathfrak{Lf}(Y)\xrightarrow{\sim}\mathfrak{Lf}(X)\\
\mathfrak{Qc}(Y)\xrightarrow{\sim}\mathfrak{Qc}(X)\\
\mathfrak{Lqc}(Y)\xrightarrow{\sim}\mathfrak{Lqc}(X)\\
\mathfrak C(Y)\xrightarrow{\sim}\mathfrak C(X)\\
\mathfrak{Lc}(Y)\xrightarrow{\sim}\mathfrak{Lc}(X).
\end{gathered}
\]
\end{proposition}

